\section{Affine returns and exhaustion by whole lines}\label{sec:lines}

Let \(C\) be a directed cycle of \(\stategraph\). Write
\begin{equation}\label{eq:cycle_length}
 \ell_C=\sum_{s\in C}h_s.
\end{equation}
The edge phase changes are their lifted lengths modulo three.
Returning to the same state therefore forces \(3\mid\ell_C\).
At any chosen base state \(s\in C\), composing the edge maps gives
\begin{equation}\label{eq:cycle_return}
 F^{\ell_C}L_s(t)=L_s(\epsilon_Ct+\beta_C),\qquad
 \epsilon_C\in\{1,-1\},\quad\beta_C\in\Z.
\end{equation}
Call \(C\) \emph{retained} if \(\epsilon_C=-1\), or if
\(\epsilon_C=1\) and \(\beta_C=0\). Moving the base state conjugates
the return map by a bijective affine parameter map, so retention is
independent of the chosen base state.

\begin{table}[ht]
\centering
\small
\begin{tabularx}{\textwidth}{@{}p{0.25\textwidth}Yp{0.27\textwidth}@{}}
\toprule
Parameter return & Decision and reason & Certified native return \\
\midrule
\(t\mapsto t\) & Retain: identity on the entire line
 & \(r_C=\ell_C/3\) \\
\(t\mapsto-t+\beta_C\) & Retain: the square is the identity
 & \(r_C=2\ell_C/3\) \\
\(t\mapsto t+\beta_C\), \(\beta_C\ne0\)
 & Do not retain: no integer parameter is periodic under translation
 & None \\
\bottomrule
\end{tabularx}
\caption{The affine-return test used in the construction. The listed
returns are not assertions about generic or exceptional least periods.}
\label{tab:returns}
\end{table}

\begin{lemma}[Periodic cycle lines]\label{lem:cycle_lines}
A nonretained cycle contains no \(F\)-periodic point on its state
lines. Every integer point on a retained cycle's state lines and all
their intermediate images is \(F\)-periodic, with a return time
\(\ell_C\) or \(2\ell_C\) as in Table~\ref{tab:returns}.
\end{lemma}

\begin{proof}
For a nonretained cycle the iterates of the parameter at the base
state are \(t+k\beta_C\), with \(\beta_C\ne0\). They are unbounded,
whereas every subsequence of a periodic orbit is bounded. For a
retained cycle the affine return has order at most two, so
\eqref{eq:cycle_return} fixes every parameter after the stated number
of lifted steps. The same return fixes every intermediate image
because powers of \(F\) commute. The edge identities transport the
assertion around the cycle.
\end{proof}

For all retained cycles form the union
\begin{equation}\label{eq:lift_line_union}
 \liftlines=
 \bigcup_{\substack{C\text{ retained}\\s\in C}}
 \ \bigcup_{0\le r<h_s}F^rL_s(\Z).
\end{equation}
Let \(\lineset\subset\Z^3\) be the phase-zero part of this union,
with the phase label omitted. The use of every intermediate image in
\eqref{eq:lift_line_union} is necessary: state lines alone need not
form an invariant set for one lifted step.

\begin{lemma}[Backward exhaustion]\label{lem:backward}
The union \(\liftlines\) satisfies
\(F(\liftlines)=\liftlines=F^{-1}(\liftlines)\).
Any point whose forward orbit reaches \(\liftlines\) already belongs
to \(\liftlines\).
\end{lemma}

\begin{proof}
Along an edge, \(F\) takes each intermediate image onto the next.
At its last image, \eqref{eq:edge_identity} and
\(\phi_s(\Z)=\Z\) show that the image is the entire next state line,
not a proper subset. As the edges run around cycles, every line image
in the union has a predecessor there. Thus
\(F(\liftlines)=\liftlines\). Bijectivity of \(F\) on the ambient
phase-labelled lattice gives the inverse-image equality. Repeatedly
applying it proves the final assertion.
\end{proof}

\begin{proposition}[Uniform line exhaustion]\label{prop:exhaustion}
Every native periodic orbit of height \(M>R\) is contained in
\(\lineset\). The set \(\lineset\) is a union of at most 27 affine
integer lines, each having a certified native return time at most 54.
\end{proposition}

\begin{proof}
By Proposition~\ref{prop:graph_cycle_entry}, the lifted orbit reaches
a directed graph cycle. Since the orbit is periodic,
Lemma~\ref{lem:cycle_lines} excludes a nonretained cycle. It therefore
reaches \(\liftlines\). Lemma~\ref{lem:backward} puts the entire
preceding orbit in the same union, including any graph path before
entry into the cycle. Its native phase-zero points lie in
\(\lineset\).

Distinct directed cycles of a graph with outdegree at most one have
disjoint state sets. A cycle of lifted length \(\ell_C\) has exactly
\(\ell_C/3\) phase-zero positions among its intermediate images,
because the phase advances by one at each lifted step. Every such
position is an affine line by Lemma~\ref{lem:edges}. Consequently the
number of phase-zero lines, before coincidences, is at most
\[
 \sum_{C\text{ retained}}\frac{\ell_C}{3}
 \le \sum_{C\text{ retained}}|C|\le27.
\]
Also \(\ell_C\le3|C|\le81\). Dividing the lifted returns in
Lemma~\ref{lem:cycle_lines} by three, using \eqref{eq:clock}, gives
\(r_C=\ell_C/3\) or \(r_C=2\ell_C/3\le54\).
\end{proof}

The assertion concerns pointwise periodicity along each line, not
setwise invariance of each separate line under \(T\). The union
\(\lineset\) is \(T\)-invariant, and its points can pass between
different lines before returning. Line coincidences only lower the
bound of 27; they create no missing family.
